\chapter{Discussion}\label{ch:discussion}

This chapter interprets the results presented in Chapter \ref{ch:results} through the lens established in Chapter \ref{ch:introduction}: the need for trustworthy and cost-aware use of \acp{llm} in \ac{se}, and the role of benchmarking as a prerequisite for reliable adoption. The discussion is organized into four parts. 

First, it interprets the empirical findings across evaluation modalities and token usage behavior. Second, it translates those findings into implications for benchmarking methodology, including task-modality alignment and evaluation robustness. Third, it derives practical guidance for deploying LLMs in SE workflows. Finally, it identifies key limitations and outlines future directions for extending this work.


\section{Interpretation of Findings}\label{sec:interpretation_findings}
\subsection{Interpreting MCQ to OSQ Conversion Suitability and Filtering Outcomes}

A central enabling step of this research is the transformation of \acp{mcq} into \acp{osq} while preserving semantic intent and domain relevance. The suitability score distribution reported in Chapter \ref{ch:results} indicates that conversion feasibility is not uniform across the source dataset. The observed split---with a substantial majority of questions meeting the inclusion threshold for OSQ conversion (845 questions, 73.9\%) and a non-trivial minority rejected (299 questions, 26.1\%), shown in Figure \ref{fig:mcq-to-osq}.

The acceptance rate suggests that domain-specific SE MCQs can be re-expressed in a textual-response form without fundamentally altering what is being tested. This supports the premise that existing benchmark assets can be expanded to richer modalities rather than rebuilt from scratch. Second, the rejected portion is not merely incidental noise; it represents structural boundaries of conversion. The rejected items include questions whose multiple-choice framing provides essential scaffolding (e.g., ambiguity resolved only by options, or questions whose discriminative power is primarily in the distractors).


\subsection{Interpreting MCQ Positional Effects and Robustness Metrics}

Position effects arise from an interaction between the model and the MCQ presentation, and can therefore be interpreted as (a) a measurement artifact, (b) a behavioral bias of the model, or (c) a combination of both.

Accordingly, the MCQ results provide two complementary forms of evidence:
\begin{enumerate}[label=(\roman*)]
    \item Outcome performance signals reflected in accuracy differences across models.
    \item Robustness signals reflected in sensitivity to controlled perturbations of the MCQ form (e.g., answer position and distractor configuration).
\end{enumerate}

The observed spread in MCQ accuracy across models (approximately 0.56 to 0.97) indicates that the dataset meaningfully separates capability levels under the evaluated prompt constraints. This supports the use of MCQs as an efficient instrument for broad capability stratification and regression testing.

However, the positional analysis motivates a more careful interpretation of what MCQ accuracy measures. If performance varies systematically with answer position or option ordering, the evaluation may partially reflect sensitivity to presentation cues rather than only domain understanding. From a benchmarking perspective, this constitutes a threat to comparability unless controlled (e.g., by randomization and variance reporting). From a model-behavior perspective, it provides a diagnostic of brittleness: models that exhibit stronger position sensitivity may be more prone to superficial heuristics that can inflate confidence in settings where the answer is not anchored by deep reasoning. This concern aligns with the trustworthiness framing in Chapter \ref{ch:introduction}, where deployment risk is driven not only by average accuracy but also by failure modes and robustness under small input changes.

Finally, \acp{mcq} trade expressiveness for efficiency. They can reward recognition and eliminate the need to articulate rationale, which is appropriate for certain SE knowledge checks (definitions, process relationships, categorical mappings). These robustness findings directly address \textbf{RQ1} and support \textbf{Contribution~2} (robustness analysis of MCQ evaluation using distractor variation) by demonstrating that positional sensitivity is a measurable and non-trivial property of MCQ-based benchmarks in the SE domain.

\subsection{Interpreting OSQ Performance}

The OSQ results introduce a different measurement regime: models must generate content, and scoring reflects rubric-based quality rather than discrete option selection. This shift inherently changes what "success" represents. A correct OSQ response can encode partial correctness, incomplete reasoning, or correct conclusions with weak justification, depending on rubric design and judging fidelity.

One important interpretive point is that OSQ scoring captures dimensions of response quality that MCQs cannot: explanatory coherence, appropriate use of SE terminology, structured reasoning, and completeness. This aligns with the motivation in Chapter \ref{ch:introduction} that textual response formats can better capture nuanced conceptual understanding for SE tasks. The discussion value of OSQs is therefore not only in ranking models, but in revealing \emph{how} they fail: omissions, shallow explanations, misapplied concepts, or fluent but incorrect reasoning.

At the same time, OSQs introduce measurement complexity. The use of kernel density estimates and the observed behavior of the density tails (an artifact of Gaussian kernels extending beyond the bounded scoring range) highlights a practical interpretive caution: OSQ score visualizations are summaries of bounded rubric outcomes, and the evaluation pipeline must ensure that interpretive claims remain grounded in the rubric domain. In other words, OSQ scoring is more expressive, but also more sensitive to judge calibration, rubric specificity, and the distributional properties of the scoring procedure.


\subsection{Interpreting the Evaluation Modality Relationship Between MCQ and OSQ }

The comparative analysis between MCQ accuracy and OSQ mean judge score provides evidence that the two modalities are related but not equivalent. The reported regression relationship (with $R^2 = 0.647$) indicates substantial shared variance, suggesting that models that perform well in MCQ tend to also perform well in OSQ, to a meaningful degree. This supports a practical interpretation: MCQ performance can serve as a coarse indicator of domain competence.

However, the fact that the relationship is not near-deterministic is equally important. Divergence between MCQ and OSQ outcomes is expected given the different cognitive demands of selection versus generation. A model may select correct answers reliably yet provide weak explanations, or conversely generate high-quality reasoning while occasionally failing brittle multiple-choice traps. In SE applications, these differences map to distinct risks. Selection-style correctness is valuable for quick checks, but generation-style competence is critical for tasks like requirements rationale, trade-space explanation, and architectural justification.


    % \item \textbf{Evaluate inter-judge agreement (correlation).}
    % Let $\rho$ denote the chosen inter-judge correlation (e.g., Spearman).
    % \begin{enumerate}[nosep]
    %     \item If $\rho \geq 0.80$: judges are strongly consistent; OSQ scores and model rankings are robust to judge choice.
    %     \item If $0.50 \leq \rho < 0.80$: judges show moderate agreement; OSQ results are usable but judge-specific effects should be acknowledged.
    %     \item If $\rho < 0.50$: judge agreement is low; treat OSQ-based conclusions with caution and highlight this as a limitation.


The cross-modality classification rubric defined in Chapter~\ref{ch:materials_methods} and applied in Chapter~\ref{ch:results} classifies the MCQ-OSQ relationship as correlated but non-interchangeable. This designation encodes two complementary findings. First, the formats produce highly correlated model rankings ($r \approx 0.80$, $\rho \approx 0.81$): a model that performs well on MCQ tends to perform well on OSQ, and vice versa. Second, MCQ systematically yields higher absolute scores than OSQ for every model evaluated (Wilcoxon $p < 0.001$, mean MCQ advantage $= 12.6$ percentage points). The effect is universal: all 19 models scored higher on MCQ, ranging from 5.3 percentage points (mistral-large) to 38.1 percentage points (llama3.2\_\_3b).

The practical implication is direct: MCQ and OSQ provide complementary but non-interchangeable evidence. MCQ captures the same relative model ordering but at a systematically inflated absolute level. OSQ reveals capability gaps that MCQ conceals. An evaluation regime that relies solely on MCQ will overestimate model readiness for tasks requiring generative reasoning, while one that relies solely on OSQ will miss the efficient screening value that MCQ provides. This directly motivates the multi-format benchmark direction laid out in Chapter~\ref{ch:introduction}: a robust SE benchmark should intentionally combine modalities to measure distinct aspects of SE-relevant competence. This cross-modality evidence directly addresses \textbf{RQ2} and provides the empirical foundation for \textbf{Contribution~3} (evaluation modality fusion as a basis for model qualification).

\subsection{Interpreting Tokenomics Factors of Quality and Verbosity}

The token distribution analysis shows that response lengths vary widely across models, from fewer than 20 tokens to more than 500 tokens for the same posed question. This provides a concrete interpretive anchor for cost-aware deployment: response verbosity is a resource consumption behavior that differs systematically across models.

The length-versus-quality analysis suggests that longer responses do not guarantee higher quality, and that quality differences persist even when controlling response length bins. This supports a core deployment interpretation: operators cannot assume that ``more tokens'' is equivalent to ``more correctness'' or ``better reasoning.'' In SE contexts where outputs may feed downstream decisions, verbosity without accuracy can: (i) provide a false sense of confidence, (ii) increase review burden, and (iii) propagate into downstream artifacts all the while preserving - if not increasing - risk.

The Pareto-style frontier visualization frames the efficiency trade space: some models achieve relatively high judged quality at moderate token usage, while others consume more tokens for comparable or lower quality. This implies that benchmark results should not only rank models by quality, but characterize efficiency regimes that inform model selection under budget constraints. The tokenomics analysis addresses \textbf{RQ3} and supports \textbf{Contribution~4} (token-efficient qualification and cost-aware response length trade-offs).


\subsubsection*{Interpretation}

% Models on the Pareto frontier achieve the best available trade-offs between cost and quality: no other model is both cheaper and more accurate. Models inside the frontier are dominated by at least one alternative that is strictly better in cost, quality, or both. The efficiency metric (e.g., cost per score point) provides a simple rank ordering of models by cost-effectiveness and enables practical recommendations such as ``cheaper models that achieve near-top performance'' versus ``premium models whose quality gains justify their cost.'' This analysis grounds benchmark results in realistic resource constraints.


Cost-effectiveness and token usage are  interpreted using the following decision logic:

\begin{enumerate}[nosep]
    \item \textbf{Examine the Pareto frontier plot.}
    \begin{enumerate}[nosep]
        \item Identify models on the upper-left boundary (high quality, low cost): these form the Pareto frontier.
        \item Models that lie strictly below and to the right of the frontier are dominated by at least one alternative and are less attractive.
    \end{enumerate}

    \item \textbf{Compare cost per 1{,}000 questions.}
    \begin{itemize}[nosep]
        \item If two models have similar quality but one has substantially lower cost, the cheaper model is preferred for large-scale deployment.
        \item If a more expensive model has noticeably higher quality, the trade-off depends on available budget and required performance.
    \end{itemize}

    \item \textbf{Evaluate the efficiency metric (e.g., cost-per-score).}
    \begin{enumerate}[nosep]
        \item Rank models by cost-per-score; lower values indicate better efficiency.
        \item Check whether highly efficient models also meet minimum quality thresholds (e.g., a target mean OSQ score).
    \end{enumerate}

    \item \textbf{Final interpretation.}
    \begin{itemize}[nosep]
        \item \emph{Operationally optimal models:} models on the Pareto frontier with low cost-per-score and acceptable quality; these are primary candidates for deployment.
        \item \emph{Premium models:} models with the highest quality but also much higher cost; preferable only when maximal performance is required and budget allows.
        \item \emph{Dominated models:} models that are simultaneously more expensive and less accurate than at least one alternative; these can be deprioritized.
    \end{itemize}
\end{enumerate}
    


\subsection{Interpreting INCOSE Task Categories}

The preceding subsections interpret aggregate, model-level performance across modalities. However, systems engineering competence is not monolithic: different INCOSE knowledge areas impose different reasoning demands and therefore different modality sensitivities. The category-level decomposition reported in Chapter~\ref{ch:results} (Section~\ref{sec:modality_mapping}) disaggregates performance across 20 INCOSE Handbook-aligned task categories and six top-level groups, enabling category-specific qualification guidance. This decomposition---summarized in Table~\ref{tab:qual_modality_class_by_incose} and visualized in Figures~\ref{fig:mcq_osq_delta_by_incose_category} and~\ref{fig:mcq_vs_osq_scatter_regions}---is arguably the most actionable result for SE practitioners, because it connects benchmark evidence directly to the task categories that govern day-to-day engineering work.

All 20 categories exhibit a positive mean MCQ--OSQ delta, confirming that the systematic MCQ inflation observed at the model level is not an artifact of aggregation but a pervasive phenomenon across the full SE knowledge taxonomy. The magnitude, however, varies substantially: from $+7.9\%$ (Architecture Definition Process) to $+19.6\%$ (Modeling and Simulation). The categories with the largest deltas---Modeling and Simulation, Modeling Frameworks and Methods, Decision Management and Tradespace Analysis, and System Safety Engineering---tend to involve synthesis-heavy, judgment-intensive knowledge where constructing a rationale is central to competence. In these areas, the MCQ format is least representative of the underlying task demand, and MCQ-only evaluation would provide the most misleading qualification evidence precisely where the stakes of generative competence are highest.

Seven categories are classified as \textit{dual-modality required} ($|\overline{\Delta}_c| \geq 0.15$): Modeling and Simulation ($+19.6\%$), Modeling Frameworks and Methods ($+18.6\%$), Decision Management, Analysis of Alternatives, and Tradespace Analysis ($+18.2\%$), System Safety Engineering ($+17.6\%$), Configuration and Information Management ($+16.5\%$), Usability and Human Systems Integration ($+15.7\%$), and Risk and Opportunity Management ($+15.4\%$). For these categories, MCQ-only evaluation carries substantial risk of overestimating model readiness for tasks that inherently require constructed justification rather than option selection. The remaining 13 categories fall in the \textit{dual-modality recommended} band ($5\% \leq |\overline{\Delta}_c| < 15\%$). Notably, no category reaches the \textit{MCQ-sufficient} classification ($|\Delta| < 5\%$ with MCQ accuracy $\geq 85\%$), which means the entire SysEngBench taxonomy exhibits meaningful modality sensitivity. This result empirically validates the Chapter~\ref{ch:introduction} premise that MCQ-only benchmarking is insufficient across the full scope of SE knowledge areas, and it provides the category-level granularity needed for task-aligned qualification decisions (\textbf{RQ2}, supporting \textbf{C1} and \textbf{C3}).

The category-level scatter plot (Figure~\ref{fig:mcq_vs_osq_scatter_regions}) reveals a tighter linear relationship between MCQ and OSQ accuracy at the category level ($R^2 = 0.811$) than at the model level ($R^2 = 0.647$). This higher explained variance suggests that when performance is averaged across models, systematic modality effects become more stable and predictable by knowledge area than by individual model. The three decision regions overlaid on the scatter---MCQ-sufficient, dual-modality recommended, and dual-modality required---provide a direct operational tool: a practitioner or benchmark designer can locate the INCOSE categories relevant to a planned deployment and immediately determine the minimum evaluation evidence required for defensible qualification.


\section{Implications for SE Benchmarking Methodologies}\label{sec:implications_benchmarking}

\subsection{Implications for Task–Modality Alignment and Systems Engineering Benchmarks}

Chapter~\ref{ch:introduction} identifies evaluation modality as a first-order design choice and argues that MCQ-only evidence is insufficient for qualifying \acp{llm} on SE tasks that require constructive reasoning. The results in Chapter~\ref{ch:results} provide empirical support for both premises and allow the implications to be stated concretely.

A primary methodological implication of this work is that benchmark validity depends on matching task type to response modality. Chapter \ref{ch:introduction} motivates this as a structured alignment problem: different SE tasks demand different evidence of competence. The Chapter \ref{ch:results} results reinforce that modality changes the measurement target. MCQs emphasize discriminative recognition, while OSQs emphasize generative reasoning and explanation quality inferred from the grading categories of technical accuracy, conceptual understanding, completeness, clarity and organization, and professional relevance.


The implication is that Systems Engineers and benchmark designers should treat modality as a first-class design decision. For concepts where the intended competence is selection among alternatives (e.g., identifying correct process steps, definitions, or categorical mappings), MCQs remain appropriate and efficient. For tasks where competence is demonstrated through justification, decomposition, or trade reasoning, OSQs provide stronger validity. 


In 15288 terms, MCQs are best aligned to measuring \textit{process literacy}—recognition of process intent, artifact mappings, and rule-based distinctions—because competence is naturally expressed as selection among alternatives and can be assessed with high reliability controls. OSQs are best aligned to measuring \textit{process execution competence}—requirements formulation, architecture and trade justification, risk rationale, and V\&V planning—because these outcomes require constructing artifacts and defensible reasoning in context. A mixed-modality benchmark therefore improves coverage across the 15288 process space by pairing efficient breadth assessment (MCQ) with high-validity depth assessment (OSQ), reducing the risk that evaluation overfits to a single response format and better represents the heterogeneous nature of SE work.


\subsection{Implications for MCQ Robustness and Reliability Checks}

Prior work has shown that MCQ performance can be sensitive to option position, distractor construction, and surface-level cues \cite{zheng_LargeLanguageModelsAre_2024, fu_mcq_reasoning_more_confident_when_wrong_2025, zhangMitigatingEasyOption2025, ghanem_DISTO_2023}.
The observed sensitivity patterns in MCQ evaluation motivate robustness practices that are uncommon in many benchmark reports. If performance varies with option ordering or distractor characteristics, then a single-form MCQ score may be an unstable and unreliable estimate of competence.

A methodological implication is that MCQ benchmarks in specialized domains should incorporate reliability checks such as:
(1) randomized answer position across evaluation runs,
(2) controlled distractor variants that probe whether models rely on superficial cues,
and (3) reporting of per-position performance rather than only aggregate accuracy.
These practices reduce the risk that benchmark results overstate model capability due to format artifacts.


\subsection{Implications for OSQ Evaluation and Judge Design  }

OSQ evaluation requires methodological discipline to remain comparable and reproducible. Rubrics should be explicit about what dimensions of quality are being measured (technical correctness, completeness, justification quality, terminology, etc.).

Because OSQ scoring often depends on automated judging, the benchmark methodology should treat the judge as part of the measurement instrument. This implies reporting judge configuration, rubric prompts, and any additional system or user prompts. Even while absolute judge reliability is imperfect, transparency allows the community to interpret scores appropriately and to compare results across studies.

The inter-judge agreement results reported in Chapter~\ref{ch:results} (Table~\ref{tab:interjudge_corr}) provide direct evidence of OSQ scoring stability. The Spearman rank correlation between the two judges (\texttt{gpt-oss\_120b} and \texttt{openai\_gpt-5-mini}) is $\rho = 0.9012$ ($p < 0.0001$, $n = 16{,}051$), indicating strong concordance in how the judges rank model responses across rubric dimensions. This level of agreement supports treating the averaged multi-judge scores as a stable measurement signal rather than an artifact of a single judge's scoring tendencies. From a methodological standpoint, the high inter-judge correlation also validates the choice of judge models: because the two judges were selected from different model families, their agreement suggests that the rubric is sufficiently specific to constrain scoring behavior and that the evaluated quality dimensions are operationally well-defined. Nevertheless, $\rho = 0.90$ is not unity, and the residual disagreement motivates reporting judge-specific results alongside the average, as done in Chapter~\ref{ch:results}, so that readers can assess whether conclusions are sensitive to judge choice.


\subsection{Implications for Reporting Standards }

A final methodological implication is that benchmarks intended to support trustworthy adoption should aim to include cross modality performance, robustness and resilience to question variance, and cost efficiency reporting rather than single modality accuracy-only comparisons on leader boards. In practice, this can be achieved by augmenting aggregate rankings with (i) modality-separated outcomes, (ii) a small robustness check where the format can introduce variants, and (iii) efficiency characterization through token distributions and quality--token trade-offs. The purpose of these additions is not to establish a new standard or replace leader boards, but to reduce ambiguity about what a score represents when results are used for procurement, workflow integration, or governance in SE contexts.


\section{Practical Applications for Systems Engineering and Insights for LLM Deployment}\label{sec:practical_applications}

\subsection{Model selection Guidance by Task Type and Risk Posture}

From a practical standpoint, the results suggest that model selection should be guided not only by average benchmark performance, but by how well a model's measured strengths align with the target task's evidence requirements and acceptable risk posture.

The table below distinguishes between (i) the \textit{recommended qualification modality}—the minimum evaluation evidence that should be satisfied before deploying a model for a task class—and (ii) \textit{modality fit}, which characterizes how well MCQ- or OSQ-style benchmarking captures the competence required for that task. Modality fit is expressed qualitatively to emphasize that benchmarking is not a binary choice: modalities provide overlapping but non-identical measurement coverage. 

Guidance is derived by identifying evaluation bundles appropriate to task criticality:

\begin{itemize}
    \item \textbf{Low-risk tasks:} MCQ-heavy screening plus targeted OSQ spot checks.
    \item \textbf{Moderate-risk tasks:} paired MCQ and OSQ evaluation on aligned subsets, with OSQ emphasized for tasks requiring explanation and rationale.
    \item \textbf{High-risk tasks:} OSQ-focused evaluation with explicit judge stability checks and robustness gating of any MCQ-derived signals.
\end{itemize}
%%%%%%%%%%%%%%%%


This aligns with the Chapter~\ref{ch:introduction} emphasis on trustworthy usage: deployment choices should be made with explicit awareness of failure modes and measurement limitations rather than relying on general-purpose claims about capability.


\begin{table}[htbp]
\centering
\caption{Notional task classes and modality guidance for benchmark-informed deployment.}
\label{tab:task_modality_guidance}
\resizebox{\linewidth}{!}{%
\begin{tabular}{p{0.7cm} p{4.0cm} p{2.2cm} p{4.2cm} p{1.5cm} p{1.5cm} p{5.6cm}}
\toprule
\textbf{\#} &
\textbf{Task class (SE use)} &
\textbf{Risk posture} &
\textbf{Recommended qualification modality} &
\textbf{MCQ fit} &
\textbf{OSQ fit} &
\textbf{Rationale / Guidance} \\
\midrule

1 &
\makecell[l]{Ideation \& coverage expansion\\(e.g., brainstorm risks/hazards,\\mitigations, questions to ask)} &
Low &
\makecell[l]{MCQ sufficient\\(OSQ optional)} &
High &
Medium &
Primary value is breadth and prompting human review. MCQ captures baseline domain recognition; OSQ can improve clarity and usefulness but is not required to qualify. \\

\addlinespace
2 &
\makecell[l]{Compression \& translation\\(summaries, rewrites,\\structured briefings)} &
Low--Medium &
\makecell[l]{OSQ preferred\\(MCQ optional)} &
Low &
High &
These tasks are inherently generative and fidelity-sensitive; OSQ better reflects coherence and completeness of the produced artifact. \\

\addlinespace
3 &
\makecell[l]{Classification \& mapping\\(taxonomy mapping,\\requirement tagging,\\process-to-standard alignment)} &
Medium &
MCQ + OSQ &
High &
High &
MCQ validates categorical recognition; OSQ validates rationale and reduces risk of ``lucky selection'' or shallow pattern matching. \\

\addlinespace
4 &
\makecell[l]{Technical interpretation \& justification\\(requirements interpretation,\\constraints/edge cases,\\V\&V rationale)} &
High &
\makecell[l]{OSQ required\\(MCQ for spot checks)} &
Medium &
High &
High-stakes reasoning requires explicit justification; OSQ surfaces reasoning quality. MCQ can be used as targeted robustness probes but should not be the qualifier alone. \\

\addlinespace
5 &
\makecell[l]{Decision support under consequence\\(trade studies, architecture\\implications, safety/mission-\\critical recommendations)} &
Very High &
\makecell[l]{OSQ + robustness controls\\+ human review} &
\makecell[c]{Low--\\Medium} &
High &
High-impact tasks require traceability and conservative governance. OSQ is necessary but remains insufficient without workflow controls and escalation rules. \\

\bottomrule
\end{tabular}
}
\vspace{0.5em}
\footnotesize
\textit{Notes.} \textbf{Recommended qualification modality} indicates the minimum evaluation evidence that should be satisfied before deploying a model for the task class. \textbf{Modality fit} is a qualitative indicator of how well MCQ- or OSQ-style benchmarking captures the competence required for the task (\textit{High} = direct measure, \textit{Medium} = partial proxy, \textit{Low} = weak proxy). In practice, high-risk uses should additionally employ workflow-level controls (e.g., human review, audit logging, escalation rules), regardless of benchmark performance.
\end{table}

\clearpage



\subsection{Benchmark-Informed SE Workflow Integration}



Benchmark results provide the evidence needed to confidently integrate LLM copilots into agentic workflows safely: models are assigned to tasks they are measured to perform well, outputs are constrained into reviewable changes, and automated validation gates reduce the likelihood that incorrect generation propagates into the model. 

A practical implication of the results is that benchmark evidence can be used to design \emph{how} LLMs are integrated into systems engineering workflows, rather than treating an LLM as a general-purpose assistant with uniform reliability. This is especially relevant in \ac{mbse}. 

Benchmark results do not merely rank models; they inform workflow roles, guardrails, and oversight practices that make copilot use auditable and dependable in an engineering context.


The INCOSE category-level results (Section~\ref{sec:modality_mapping}) sharpen this perspective by providing category-specific evidence for role assignment. Categories classified as \textit{dual-modality required}---such as System Safety Engineering, Decision Management, and Modeling and Simulation---exhibit large MCQ--OSQ deltas that indicate MCQ-only screening is unreliable for these topics. In a benchmark-informed workflow, tasks falling within these categories should be routed to models that have demonstrated strong OSQ performance in the relevant domain area, and their outputs should be subject to rubric-caliber review rather than binary correctness checks. Conversely, tasks in categories with smaller modality gaps (e.g., Architecture Definition Process, Identifying Stakeholder Needs) may tolerate MCQ-screened models with lighter oversight, reducing the overall cost of quality assurance.

Beyond model selection, the results motivate augmentation strategies that address identified capability gaps. Models that perform well on MCQ but poorly on OSQ in specific categories may benefit from retrieval-augmented generation using domain-specific knowledge bases (e.g., INCOSE Handbook sections, organizational standards, or project repositories) to provide grounding context that improves generative response quality. Similarly, fine-tuning on domain-specific corpora may close modality gaps for models that exhibit the requisite base capability. The benchmark results provide the diagnostic evidence needed to target such interventions: rather than applying augmentation uniformly, practitioners can prioritize categories where the MCQ--OSQ delta is largest and where the task risk posture demands generative competence.

In \ac{mbse} contexts, benchmark-informed role assignment maps naturally to the model-based workflow. Models assigned to structured validation roles (e.g., checking requirement completeness against a template, verifying naming conventions, or performing taxonomy alignment) can operate with MCQ-calibrated confidence. Models assigned to constructive roles (e.g., drafting trade rationale, generating architecture justification, or composing risk narratives for review) should be qualified using OSQ evidence in the relevant INCOSE categories. This mapping treats the benchmark not as a one-time model ranking exercise but as a persistent qualification instrument that informs workflow configuration, role boundaries, and oversight intensity.

The Chapter~\ref{ch:results} findings indicate that performance varies by evaluation modality and that response behavior (including verbosity and variability) differs across models. This motivates treating copilot selection as a role assignment problem. In practice, a benchmark-informed workflow assigns models to tasks aligned with demonstrated strengths. For example, models that perform consistently on structured selection-style evaluations can be used for fast conformance checks, taxonomy mapping, and structured validation prompts, while models that score strongly on open-response evaluation can be used for drafting rationale, explanations, and structured narrative artifacts. Conversely, models exhibiting unfavorable quality--token trade-offs or high variability may be restricted to low-risk tasks, short-form assistance, or preliminary drafts to reduce review burden and operational cost.

This role assignment perspective supports the Chapter~\ref{ch:introduction} motivation for evidence-based adoption: instead of asking whether a model is ``good'' in general, the workflow asks whether a model is suitable for a specific role under defined constraints and oversight.





\subsection{Budgeting and Token Constraints}

The tokenomics results translate directly into deployment policy. Since verbosity varies widely across models, organizations can manage cost and review burden by enforcing response budgets and adopting structured prompting patterns. Examples include:
(i) requesting bounded-length responses with required sections,
(ii) using iterative refinement (short first-pass, expanded only when needed),
and (iii) separating reasoning generation from final answer formatting.

Crucially, the length-versus-quality findings indicate that token limits do not automatically degrade quality if prompts are designed to elicit concise but complete reasoning. This supports the Chapter~\ref{ch:introduction} motivation that cost-aware adoption is feasible when guided by evidence rather than assumption.



\section{Limitations and Future Work}\label{sec:limits_and_future_work}

This section discusses the limitations of the research and outlines future research directions. Limitations are organized along the following dimensions: scope, dataset construction, execution constraints, model behavior, and evaluation reliability. Future work synthesizes research findings with limitations into a forward-looking, benchmark-informed systems engineering workflow.

\subsection{Scope and Temporal Constraints}

This dissertation is scoped to evaluate \acp{llm} for systems engineering tasks through comparative \ac{mcq} and \ac{osq} benchmarking, including hybrid approaches and tokenomics analysis. The work is anchored in foundational systems engineering concepts drawn from the INCOSE Handbook and related activities such as requirements analysis, conceptual design, risk management, and \ac{mbse}. The overarching goal is to make SysEngBench a community-driven, extensible benchmarking resource for systems engineering, supplying a methodological template for both its ongoing development and for those creating private or derivative benchmarks.

The research is deliberately bounded in both concept and method. It examines model performance and evaluation modalities with benchmark modalities are confined to \ac{mcq} and \ac{osq} formats and their hybrid combinations; multi-modal inputs and agentic workflows remain outside the present scope. From a data perspective, SysEngBench relies available systems engineering sources and does not incorporate proprietary industrial or government datasets. Temporally, the evaluation reflects the 2024 to late-2025 \ac{llm} landscape, recognizing that future models may shift absolute performance while leaving the analytical framework relevant.

\subsection{Dataset and Taxonomy Limitations}

Several dataset and taxonomy limitations motivate future extensions to SysEngBench and the evaluation of \ac{ai} models for Systems Engineers.

\subsubsection{Lack of Difficulty and Cognitive Depth Annotation}

SysEngBench lacks explicit annotations for question difficulty (e.g., easy, medium, hard) and for cognitive depth, such as Bloom's taxonomy \cite{blooms_taxonomy_1956}. This gap was partially addressed in the MCQ-to-OSQ conversion prompt (\ref{app:converter_prompt_template}) to preserve data for future work; however, this coverage applies only to converted items. MCQs that were not converted were neither labeled nor human reviewed.

As a result, the benchmark cannot report performance conditioned on difficulty strata, nor can it isolate whether observed performance differences are driven by genuine capability variation or by uneven mixtures of question difficulty across categories. Without difficulty or cognitive depth labels, SysEngBench may under-characterize where models provide meaningful systems engineering value.

Prior work has demonstrated the value of structuring evaluation benchmarks by Bloom levels for deeper model assessment \cite{Liangtai_Sun_2023_SciEval_Benchmark}. Extending SysEngBench to include Bloom-level labeling would enable more diagnostic analyses, such as distinguishing recognition-dominated MCQ performance from OSQ responses that reflect deeper reasoning and synthesis.

\subsubsection{Category Imbalance and Coverage Skew}

Category representation remains uneven across the benchmark. Some systems engineering concept categories contain substantially more items than others, which can bias aggregate conclusions toward heavily represented topics. This imbalance is particularly consequential when categories differ in inherent difficulty. Future SysEngBench releases should expand sparse categories and publish per-category performance summaries to reduce the risk of overgeneralization from aggregate results.

\subsection{Evaluation and Execution Constraints}

Several limitations arose from practical considerations during model execution rather than from the benchmark design itself.

\subsubsection{Token Budget Constraints}

The initial evaluation configuration enforced a maximum response length of 10,000 tokens. Under this constraint, several thinking-oriented models frequently exhausted the token budget before producing a final response suitable for scoring. To distinguish truncation artifacts from genuinely non-terminating behavior, the maximum response budget was decreased to 500 tokens for a targeted subset of models exhibiting repeated truncations or incomplete outputs.

Despite this adjustment, some models continued to produce operationally unusable outputs, such as indefinite generation or prolonged internal reasoning without a terminal answer. Models that consistently provided these outputs were omitted from quantitative analyses to preserve comparability and to avoid scoring artifacts driven by incomplete responses rather than task competence.

\subsubsection{Inference Timeouts and Service Reliability}

A small number of evaluation responses within less than 4 runs experienced timeouts that were attributable to operational conditions rather than to specific model behaviors. These timeouts are attributed to high concurrency (multi-threaded execution with a large number of simultaneous requests) or due to inference-service-level effects in OpenRouter (e.g., transient throttling, queuing, or upstream routing latency) during peak load. Fewer than 27 such occurrences were observed across all models and they were clustered during high-parallelism intervals. Timeout responses were deterministically re-run with temperature set to zero and completed successfully, indicating that the failures were transient execution artifacts rather than systematic model limitations.

\subsubsection{Idempotence and Deterministic Re-Execution}

The evaluation pipeline and associated code were designed to be idempotent. Aside from transient networking failures, repeated executions under identical configurations yielded consistent results. This supports the reproducibility of reported findings and isolates observed variability to model behavior rather than to execution nondeterminism.

\subsection{Model Behavior and Reasoning Control Limitations}

\subsubsection{Thinking-Model Failure Modes}

A subset of model runs were omitted due to response-format incompatibilities or pathological generation behaviors. Although the \ac{mcq} and \ac{osq} pipelines were designed to support broad model coverage, the empirical results reflect only those models whose outputs could be parsed and scored reliably under the study conditions. This introduces a selection effect favoring models that both produce responses within expected structural bounds and remain stable under the prompting and token limits used.

For models that expose explicit reasoning traces, internal "thinking" content counted toward the response token budget. In cases where a model remained in a reasoning phase without emitting a final answer, no scorable output was produced, resulting in low or null scores. These behaviors were initially observed when running a test batch of fewer than 50 evaluation items per affected model, where reliability in producing bounded, terminal responses remained insufficient despite attempts at mitigation.

\subsubsection{Incompatibilities with Reasoning Control Flags}

Additional limitations arose from heterogeneous support for reasoning-control parameters across inference stacks, particularly when running models with Ollama. Attempts to suppress or hide reasoning traces varied across model choices, including Phi-4, including:

\begin{itemize}
    \item \textbf{Standard settings:} models emitting reasoning prior to final answers broke naive MCQ parsing by introducing early option tokens during explanation.
    \item \textbf{\texttt{think=false}:} some models continued to emit reasoning traces, indicating incomplete compatibility with inference controls.
    \item \textbf{\texttt{hidethinking=true}:} suppressing reasoning traces sometimes resulted in empty-string outputs that were not scorable if response went over the token limit.
\end{itemize}


\subsubsection{Selection Effects Introduced by Model Filtering}

The decision was made to exclude open-source thinking models hosted in environments where reasoning could not be reliably controlled or parsed. The focus of this dissertation is meta-evaluation methodology for systems engineering benchmarks rather than benchmarking reasoning trace behavior itself.

Reasoning-capable proprietary API models (e.g., GPT-5) typically withhold raw chain-of-thought by default, exposing only final answers and, optionally, reasoning summaries \cite{openai_ReasoningModels_2025}. This reduces exposure to leakage-style failure modes (e.g., revealing unsafe intermediate deliberation) and encouraged treating deliberation as non-observable across evaluated systems.


\subsection{Scoring, Parsing, and Judge Reliability}

\subsubsection{MCQ Parsing Fragility}

In MCQ settings, explicit reasoning traces occasionally caused incorrect answer extraction when models enumerated options during explanation. Naive regular-expression approaches could capture an early option token rather than the final selection. Correcting this would require more robust parsing strategies (e.g., explicit final-answer anchoring) or manual intervention, both of which were deemed outside the scope of the framework used in this study.

\subsubsection{OSQ Non-Termination and Empty Outputs}

During OSQ evaluation runs, some models failed to produce judge-able responses due to indefinite generation, token exhaustion, or suppressed outputs yielding empty strings. These behaviors prevented reliable scoring and motivated model exclusion of affected outputs from aggregate analyses.

\subsubsection{Implications for Automated Benchmarking Pipelines}

These parsing and termination issues highlight limitations in current automated benchmarking pipelines when applied to heterogeneous model behaviors. They reinforce the importance of explicit output contracts (e.g. JSON), termination guarantees, and parsing robustness in future benchmark designs.


\subsection{Future Directions}

This dissertation establishes a foundation for qualification-oriented evaluation of large language models in systems engineering and the long-term trajectory extends toward a systematic qualification approach for AI usage within SE.

\paragraph{Cognitive-Difficulty as a Structural Benchmark Dimension.}
Future versions of SysEngBench should incorporate mandatory cognitive-depth labeling (e.g., Bloom-level or an equivalent engineering reasoning taxonomy) and report performance conditioned on difficulty rather than aggregate accuracy alone. Difficulty-aware reporting would enable qualification decisions aligned to task complexity and help identify capability cliffs that remain obscured in overall averages.

\paragraph{Multi-modal Artifact Integration.}
Systems engineering artifacts are not purely textual, and future benchmarking should include diagram interpretation, text–model alignment, and cross-artifact consistency tasks reflective of the multi-modal nature of \ac{se}. Evaluating vision-language models and multi-modal pipelines on these representations would provide increased qualification evidence more directly aligned to model-based engineering workflows.

\paragraph{Human–AI Performance Comparison.}
Future research could systematically compare strengths, limitations, and failure modes of humans and LLMs across representative engineering tasks, establishing empirically grounded competency metrics for LLMs relative to that of human, systems engineering experts. Such studies, likely requiring Institutional Review Board oversight, would support mapping safe operating envelopes and defining regions where human oversight remains essential.

\paragraph{Structured Multi-Agent.}
There are many circumstances within \ac{se} that could leverage multi-agent deliberation mechanisms similar to that of domain subject matter experts adjudication. Systematically studying consensus formation, ensemble disagreement, and adversarial critique could provide additional information for systems engineers to make informed trade space decisions.

\paragraph{Cost-Aware AI Qualification.}
Building on the token-efficiency analysis presented in this dissertation, future research should formalize economic models that capture inference cost, robustness replication overhead, and scoring expenses. Such models would support trade studies that determine how much evaluation evidence is sufficient for a given risk posture and budget constraint.

